% Requires amsmath, amssymb, and amsthm. No priority claim is made.
\newtheorem{auditproposition}{Proposition}[section]
\section{Identifying the contribution of structural metadata}
\label{sec:identifiability}

An external scaling policy may use a row's coverage status, a kernel selected by its
role, and the ownership of its columns. These are distinct inputs. A comparison
identifies the contribution of block ownership only when coverage, the row-role
assignment, and the kernels remain fixed. Matching the local kernel alone does not
meet this condition. This distinction has a particularly strong consequence for
the post-normalization block statistic used by the original policy.

\subsection{Collapse of post-normalization block centers}

For each nonzero local row $r$, write $m_r$ and $M_r$ for its smallest and largest
nonzero coefficient magnitudes. The coefficient geometric-mean kernel applies
$d_r=(m_rM_r)^{-1/2}$. Let $\widetilde m_r=d_rm_r$ and
$\widetilde M_r=d_rM_r$ denote the exported extrema. For a nonempty set $I$ of local
rows, define
\[
 s_I=\sqrt{\left(\min_{r\in I}\widetilde m_r\right)
                 \left(\max_{r\in I}\widetilde M_r\right)}.
\]
This definition uses all coefficients, including the row extrema.

\begin{auditproposition}[Exact collapse and the effect of a skip band]
\label{prop:audit-collapse}
If every row receives its geometric-mean factor, then $s_I=1$ for every
nonempty grouping $I$ of local rows. More generally, suppose the only safeguard is
to replace $d_r$ by $1$ when its candidate belongs to $(1/u,u)$, where $u\geq1$.
Then $s_I\in[1/u,u]$, independently of the original coefficient scales and of the
grouping of the rows.
\end{auditproposition}
\begin{proof}
Under exact normalization,
$\widetilde m_r=\sqrt{m_r/M_r}=1/\widetilde M_r$. Consequently
$\min_r\widetilde m_r=1/\max_r\widetilde M_r$, proving the first statement.
For the second, let $c_r=\sqrt{\widetilde m_r\widetilde M_r}$. An accepted row has
$c_r=1$; a skipped row has $c_r=\sqrt{m_rM_r}\in(1/u,u)$. Thus
$u^{-2}\leq\widetilde m_r\widetilde M_r\leq u^2$ for every row. Choose $p$ attaining
the maximum exported row magnitude and $q$ attaining the minimum. Then
\[
 u^{-2}\leq\widetilde m_q\widetilde M_q
 \leq\widetilde m_q\widetilde M_p
 \leq\widetilde m_p\widetilde M_p\leq u^2.
\]
Taking square roots gives the result.
\end{proof}

For a coupling row $a_r$, let $\beta(j)$ assign column $j$ to a local block, and
define the block-weighted candidate and its ownership-blind counterpart by
\[
 \gamma_r^{\mathrm{block}}
 =\left(\sum_j a_{rj}^2/s_{\beta(j)}^2\right)^{-1/2},
 \qquad
 \gamma_r^{\mathrm{L2}}=\left(\sum_j a_{rj}^2\right)^{-1/2}.
\]
Blockless columns may use an arithmetic mean of valid block scales or the unit
fallback; both conventions preserve the interval in the preceding proposition.

\begin{auditproposition}[Loss of ownership information]
\label{prop:audit-ownership}
Under exact local normalization and with no coefficient filtering,
$\gamma_r^{\mathrm{block}}=\gamma_r^{\mathrm{L2}}$ for every coupling row. Thus
changing the block ownership map cannot change the candidate export. Applying
the same coupling safeguards to the identical candidates preserves equality.
Under the skip-band assumptions of Proposition~\ref{prop:audit-collapse},
\[
 u^{-1}\leq\frac{\gamma_r^{\mathrm{block}}}
                       {\gamma_r^{\mathrm{L2}}}\leq u.
\]
Before coupling safeguards, the condition numbers of the two exports differ by
at most a factor $u^2$, on their common nonzero singular spectrum.
\end{auditproposition}
\begin{proof}
The equality follows by substituting $s_i=1$. In the second case,
\[
 u^{-2}\|a_r\|_2^2\leq\sum_j a_{rj}^2/s_{\beta(j)}^2
 \leq u^2\|a_r\|_2^2,
\]
which gives the factor bound. Write the block-based export as $EA_{\mathrm{L2}}$,
where $E$ has unit entries on local rows and entries in $[1/u,u]$ on coupling
rows. Multiplication by $E$ changes each singular value by a factor between
$1/u$ and $u$ and preserves rank. Applying these inequalities to the largest and
smallest positive singular values in both directions proves the condition-number
bound.
\end{proof}

Coefficient-window vetoes and filtering must be audited separately: they are
outside the hypotheses. Even exact row normalization does not preserve the
reciprocity argument if the block summary discards extrema. For example, a row
with magnitudes $(10^{-8},3,10^8)$ is geometrically centered, but restricting the
summary to $(10^{-6},10^6)$ produces $s_I=3$. Such a nonunit statistic arises
from the filter. It cannot, by itself, be interpreted as retained physical scale.

\subsection{What a matched comparison identifies}

With equal coverage and equal local kernels, two exports agree on their local
rows. Their differences are confined to the coupling rows. This elementary
localization result does not establish that the difference is due to block
ownership: the coupling kernels may differ too. An appropriate ownership control
keeps local geometric-mean normalization and coupling Euclidean normalization,
but replaces all block scales by one. A separate role control compares this
hybrid policy with Euclidean normalization on every row. Comparing only with
geometric-mean normalization on every row mixes role-dependent kernel selection
with ownership information.

Inactivity of one policy's coupling stage does not imply equality with a flat
policy. Consider local rows $x_1\leq1$, $x_2\leq1$, and the coupling inequality
$x_1+10^{-2}x_2\leq1$. Both local blocks have scale one. The Euclidean coupling
candidate is $(1+10^{-4})^{-1/2}$ and is skipped by the band $(1/2,2)$. The flat
geometric-mean candidate is $10$ and is applied. Both policies cover all rows,
yet their coupling exports differ. Equality requires equality of the
\emph{actually applied} factors, rather than inactivity of just one policy.

\subsection{A valid explanation of the stylized spectral slope}

Take $B\geq2$ local singleton rows with positive coefficients $v_i=10^{e_i}$,
where $\min_i e_i=-R$ and $\max_i e_i=R$, and one coupling row $v^\top$.
Exact local normalization produces $I_B$ and unit post-local block centers. A
whole-row coupling factor $\gamma$ gives
\[
 A_\gamma=\begin{bmatrix}I_B\\\gamma v^\top\end{bmatrix},
 \qquad A_\gamma^\top A_\gamma=I_B+\gamma^2vv^\top,
 \qquad
 \kappa_2(A_\gamma)=\sqrt{1+\gamma^2\|v\|_2^2}.
\]
Flat geometric-mean normalization selects $\gamma=1$, whereas both the
post-local block rule and the hybrid Euclidean control select
$\gamma=1/\|v\|_2$. Their condition numbers are therefore
\[
 \kappa_2(A_{\mathrm{GM}})=\sqrt{1+\sum_i10^{2e_i}},
 \qquad
 \kappa_2(A_{\mathrm{block}})=\kappa_2(A_{\mathrm{hybrid}})=\sqrt2.
\]
For fixed $B$, the ratio is of order $10^{-R}$. The slope is explained by the
coupling kernel and survives removal of the ownership map. Every coupling
coefficient is multiplied by the \emph{same} $\gamma$, and all within-row ratios
are preserved. Replacing $v_i$ individually by $v_i/10^{e_i}$ would change the
constraint and is not a row-scaling operation.

The matrix $A_\gamma$ also illustrates a limitation of this spectral endpoint:
letting $\gamma$ tend to zero drives $\kappa_2$ to one while allowing an
arbitrarily large original-row residual under a fixed exported residual budget.
Thus minimizing this rectangular-matrix condition number alone does not enforce
the intended numerical accuracy of the coupling constraint.

\subsection{Exact scalar optimization in the triangular model}

The idealized triangular model also admits a useful closed form that requires
no ownership map after local whitening. It describes a common coupling factor,
not the per-row optimum for a general generated MIP.

\begin{auditproposition}[Optimal common coupling factor after local whitening]
\label{prop:audit-triangular-optimum}
Let $R\in\mathbb R^{m\times n}$, with $m,n\geq1$, and set $\rho=\|R\|_2$. For
\[
 C_\gamma=\begin{bmatrix}I_n&0\\\gamma R&\gamma I_m\end{bmatrix},
 \qquad\gamma>0,
\]
the unique factor minimizing $\kappa_2(C_\gamma)$ is
\[
 \gamma_*=(1+\rho^2)^{-1/2},
 \qquad
 \min_{\gamma>0}\kappa_2(C_\gamma)=\sqrt{1+\rho^2}+\rho.
\]
With an admissible interval $0<\ell\leq\gamma\leq h$, the optimum is
$\min\{h,\max\{\ell,\gamma_*\}\}$.
\end{auditproposition}
\begin{proof}
An SVD of $R$ and orthogonal transformations reduce the nontrivial spectrum to
blocks $\left[\begin{smallmatrix}1&0\\\gamma\sigma&\gamma\end{smallmatrix}\right]$,
one for each singular value $\sigma$ of $R$, with possible additional singular
values $1$ or $\gamma$. Each block's squared singular values have product
$\gamma^2$ and sum $1+\gamma^2(1+\sigma^2)$. Their largest value increases and
smallest value decreases with $\sigma$; the block at $\sigma=\rho$ therefore
determines both extrema, including any additional values. If $k$ is their
singular-value ratio, then
\[
 k+k^{-1}=\frac{1+\gamma^2(1+\rho^2)}{\gamma}
 =\gamma^{-1}+\gamma(1+\rho^2).
\]
The left-hand side is increasing for $k\geq1$. The right-hand side has a unique
minimum at $\gamma_*=(1+\rho^2)^{-1/2}$, decreases before that point, and increases
after it. At the minimum, $k+k^{-1}=2\sqrt{1+\rho^2}$, whose root $k\geq1$ is
$\sqrt{1+\rho^2}+\rho$. The interval result follows from the same monotonicity.
For $\rho=0$, the argument reduces to balancing the diagonal entries at one.
\end{proof}

For a single coupling row, $\gamma_*$ normalizes that complete row's Euclidean
norm, including the aggregate-variable coefficient. This provides an exact
kernel explanation of the idealized conditioning improvement. It is not a
guarantee about LP bases, presolve transformations, or mixed-integer runtime.
